\chapter{Resultados Detallados de Análisis}

\section{Anexo A: Diagrama Unifilar Simplificado del Sistema}

\subsection{Descripción General}

El sistema eléctrico del proyecto solar fotovoltaico MAS X MENOS está compuesto por los siguientes elementos principales:

\begin{enumerate}
    \item \textbf{Paneles Solares:} 240 paneles JINKO monocristalinos de 590 W cada uno, total 141.6 kWp
    
    \item \textbf{Inversores:} 2 inversores SOLIS 60K-LV-5G de 60 kW cada uno, total 120 kW en AC
    
    \item \textbf{Transformador de Conexión:} 150 kVA, 13.2 kV / 220 V, grupo Dyn5
    
    \item \textbf{Protecciones en Inversor:} Breaker trifásico 3 × (140-200) A
    
    \item \textbf{Protecciones en Acometida:} Breaker trifásico 3 × (350-500) A
    
    \item \textbf{Punto de Conexión:} Nodo 3272966, red de ESSA a 13.8 kV
\end{enumerate}

\subsection{Diagrama de Bloques}

\begin{center}
    \begin{tabular}{c}
    \toprule
    \textbf{COMPONENTES DEL SISTEMA} \\
    \midrule
    \\
    Paneles Fotovoltaicos \\
    (240 × 590 W = 141.6 kWp) \\
    \textbar \\
    \textbar \\
    Inversores SOLIS \\
    (2 × 60 kW = 120 kW AC) \\
    \textbar \\
    \textbar \\
    Protección de Inversor \\
    (3 × 140-200 A) \\
    \textbar \\
    \textbar \\
    Protección de Acometida \\
    (3 × 350-500 A) \\
    \textbar \\
    \textbar \\
    Transformador de Conexión \\
    (150 kVA, 13.2 kV/220 V, Dyn5) \\
    \textbar \\
    \textbar \\
    Red de ESSA \\
    (Nodo 3272966, 13.8 kV) \\
    \\
    \bottomrule
    \end{tabular}
\end{center}

\section{Tabla de Especificaciones Técnicas Resumen}

\begin{table}[H]
    \centering
    \caption{Especificaciones técnicas generales del proyecto}
    \begin{tabular}{ll}
    \toprule
    \textbf{Parámetro} & \textbf{Valor} \\
    \midrule
    Ubicación & Bucaramanga, Santander \\
    Latitud & 7.12621563° \\
    Longitud & -73.11227598° \\
    Capacidad instalada DC & 141.6 kWp \\
    Capacidad instalada AC & 120 kW \\
    Tecnología & Paneles monocristalinos JINKO \\
    Número de paneles & 240 \\
    Potencia por panel & 590 W \\
    Eficiencia de paneles & 22.84\% \\
    Número de inversores & 2 \\
    Potencia por inversor & 60 kW \\
    Factor de potencia & 0.99 \\
    Tensión de salida & 220 VAC, 60 Hz \\
    Relación DC/AC & 1.18 \\
    Transformador & 150 kVA, 13.2/0.22 kV \\
    Punto de conexión & Nodo 3272966 \\
    Nivel de tensión conexión & 13.8 kV \\
    Año de entrada en operación & 2025 \\
    \bottomrule
    \end{tabular}
\end{table}

\section{Matriz de Evaluación de Impacto}

\begin{table}[H]
    \centering
    \caption{Evaluación de impacto en parámetros del sistema}
    \begin{tabular}{lcccc}
    \toprule
    \textbf{Parámetro} & \textbf{Límite} & \textbf{Máximo} & \textbf{Cumple} & \textbf{Impacto} \\
    \midrule
    Tensión (p.u.) & 0.90-1.10 & 0.9953 & ✓ & Mínimo \\
    Cargabilidad líneas (\%) & ≤ 100 & 4.57 & ✓ & Mínimo \\
    Cargabilidad trafo (\%) & ≤ 100 & 55.83 & ✓ & Bajo \\
    Variación CC (kA) & < 1\% & 0.00\% & ✓ & Nulo \\
    Factor de potencia & ≥ 0.90 & 0.99 & ✓ & Positivo \\
    Pérdidas técnicas (kW) & Reducción & -0.34 & ✓ & Positivo \\
    \bottomrule
    \end{tabular}
\end{table}

\section{Cronograma de Validación Recomendado}

\begin{table}[H]
    \centering
    \caption{Actividades de validación post-conexión}
    \begin{tabular}{llcl}
    \toprule
    \textbf{Fase} & \textbf{Actividad} & \textbf{Plazo} & \textbf{Responsable} \\
    \midrule
    Pre-operación & Inspección de instalaciones & 2 semanas & COPOWER + ESSA \\
    Pre-operación & Pruebas de protecciones & 2 semanas & COPOWER + ESSA \\
    Operación & Mediciones de calidad & 1 mes & COPOWER + ESSA \\
    Operación & Validación de tensiones & 3 meses & ESSA \\
    Operación & Reporte de desempeño & Trimestral & COPOWER \\
    \bottomrule
    \end{tabular}
\end{table}

\section{Documentos de Referencia}

Durante el desarrollo de este estudio se utilizaron los siguientes documentos:

\begin{itemize}
    \item Resolución CREG 030 de 2018 - Conexión de generadores
    \item Resolución CREG 025 de 1995 - Código de Redes
    \item Resolución CREG 174 de 2021 - Generación distribuida
    \item Acuerdo 1322 del CNO - Protecciones de generadores
    \item Norma IEC 60909 - Cálculo de cortocircuito
    \item Norma IEEE 242 - Protecciones de sistemas
    \item Datos de Circuito 10502 - ESSA
    \item Ficha técnica inversor SOLIS 60K-LV-5G
    \item Ficha técnica paneles JINKO Monocristalinos
\end{itemize}

\section{Contacto y Consultas}

Para consultas técnicas adicionales respecto a este estudio, contactar a:

\begin{center}
    \textbf{COPOWER LTDA} \\
    Especialistas en Energías Renovables \\
    Bucaramanga, Santander - Colombia \\
    Julio de 2026
\end{center}
